\section{Industry Adoption by Vertical}
\label{sec:adoption}

This section addresses Contribution C3
(\S\ref{sec:intro:contributions}): mapping inter-agent
protocol adoption across four verticals. The verticals
are enterprise SaaS, agentic commerce and payments,
healthcare, and finance. We name production deployments
rather than treating protocols in the abstract.
Section~\ref{sec:adoption:saas}
covers the enterprise SaaS cluster where A2A and MCP have
both achieved generally-available status across multiple
named vendors. Section~\ref{sec:adoption:commerce} treats
the agentic commerce stack, where the deployment pattern is
inverted: a small number of major payment networks rather
than many SaaS vendors. Section~\ref{sec:adoption:healthcare}
and Section~\ref{sec:adoption:finance} treat the
heavily-regulated verticals where public disclosure is
sparser; we report what is publicly known and flag the
discovery problem the survey faces in mapping these
sectors. Table~\ref{tab:industry-adoption} consolidates the
findings.

We are weighing the table according to source tier
(\S\ref{sec:intro:notwhat}, R4): every status claim
(\emph{GA} / \emph{Pilot} / \emph{Announced}) is sourced
from vendor blogs, press releases, or developer
documentation (Tier-4 in the survey's tier ranking).
Each claim is used only for the binary ``\emph{X} announced
deploying \emph{Y} on date \emph{Z}'' assertion. The table
makes no behavioural claims about the deployments
themselves; the rubric in
\S\ref{sec:rubric} is the appropriate instrument for
those. The behavioural-versus-deployment-status
distinction is critical: a vendor's claim that their
agent is ``A2A-compliant'' is a status claim;
the agent's actual rubric profile (which axes it
implements, which it leaves out) is a separate question
that vendor announcements rarely answer.

% Generated by scripts/gen_table4.py: do NOT hand-edit.
% Source: data/industry-adoption.yaml.

\begin{table*}[t]
  \centering
  \caption{Industry adoption by vertical at the 2026-Q1 cutoff. Cells list named production-or-pilot deployments with status: \textit{GA} (generally available), \textit{Pilot}, or \textit{Ann.}\ (announced). Status claims are Tier-4 (vendor blog) sourced; behavioural claims about deployments are not made in this table. The healthcare and finance verticals are under-populated because public disclosure is less common in those sectors.}
  \label{tab:industry-adoption}
  \footnotesize
  \begin{tabular}{p{2.3cm}p{2.6cm}p{2.6cm}p{2.6cm}p{2.6cm}p{2.6cm}}
    \toprule
    Vertical & A2A & MCP & ACP & AGNTCY & Commerce stack \\
    \midrule
    Enterprise SaaS & \scriptsize Salesforce Agentforce (GA), \linebreak ServiceNow AI Agents (GA), \linebreak SAP Joule (Pilot), \linebreak Atlassian Rovo (GA), \linebreak Microsoft Copilot Studio (GA), \linebreak Box AI Agents (GA), \linebreak Workday Agent System (Pilot) & \scriptsize GitHub Copilot (GA), \linebreak Anthropic Claude Code (GA), \linebreak Cursor IDE (GA), \linebreak Replit Agent (GA), \linebreak Sourcegraph Cody (GA), \linebreak JetBrains AI Assistant (GA) & \scriptsize IBM watsonx Orchestrate (GA), \linebreak IBM BeeAI internal (Pilot) & \scriptsize Cisco Webex AI Agent (Pilot) & n/a \\[6pt]
    Commerce \& payments & \scriptsize Shopify (A2A merchant agents) (Pilot) & \scriptsize Stripe MCP (GA) & n/a & n/a & \scriptsize PayPal (AP2) (GA), \linebreak Visa TAP (GA), \linebreak Mastercard Agent Pay (Pilot), \linebreak American Express (Ann.), \linebreak ERC-8004 reference impls (Pilot) \\[6pt]
    Healthcare & \scriptsize Epic agent connectivity (Pilot), \linebreak Cedars-Sinai + Anthropic (Pilot) & \scriptsize Hippocratic AI (internal) (GA), \linebreak Nuance DAX Copilot (GA) & n/a & n/a & n/a \\[6pt]
    Finance & \scriptsize Bloomberg AI (selected feeds) (Pilot) & \scriptsize BlackRock Aladdin AI (Pilot), \linebreak JPMorgan IndexGPT family (Pilot) & n/a & n/a & \scriptsize DTCC (settlement agent pilots) (Pilot) \\[6pt]
    \bottomrule
  \end{tabular}
\end{table*}


\subsection{Enterprise SaaS}
\label{sec:adoption:saas}

The enterprise SaaS vertical is the densest deployment
cluster at the cutoff and the one where the
A2A-versus-MCP composition pattern
(\S\ref{sec:taxonomy:complementary}) is most visible in
production. All status claims in this subsection are
sourced from public vendor announcements catalogued in
\texttt{data/industry-adoption.yaml} of the companion
repository~\cite{vendor_announcements_2026}; we cite
specifications directly only where behavioural rather
than status claims are made.

\paragraph{A2A deployments.}
Salesforce Agentforce, ServiceNow AI Agents, Microsoft
Copilot Studio, Atlassian Rovo, and Box AI Agents are at
generally-available status with A2A integration as of the
cutoff. SAP Joule and Workday Agent System are in pilot
status. The deployment pattern is consistent across the
cluster: each vendor exposes its SaaS-native agents
behind A2A AgentCards, allowing third-party orchestrators
(or competing vendor's agents) to invoke them as remote
agents over A2A. The interoperability claim is bounded:
A2A's discovery convention permits cross-vendor calls in
principle. But Section~\ref{sec:semantic:fragmentation}'s
schema fragmentation risk applies in full. A
billing query that succeeds against Salesforce
Agentforce may fail or silently mis-execute against
ServiceNow's billing agent. The two vendors
authored their Skill descriptions independently.

\paragraph{MCP deployments.}
MCP's enterprise SaaS deployment is in developer-tooling
rather than business-application territory: GitHub Copilot,
Anthropic Claude Code, Cursor IDE, Replit Agent,
Sourcegraph Cody, and JetBrains AI Assistant all ship MCP
support at GA. The cluster reflects MCP's tool-binding
positioning: developer-tooling vendors expose
code-search, repository operations, build-system
invocation, and similar capabilities through MCP servers
that any of the surveyed clients can consume.
The deployment population is more homogeneous than
the A2A cluster. Developer-tooling vendors have
historically converged on shared interfaces faster than
business-application vendors. MCP's lower authoring
overhead (a single JSON Schema per tool) accelerates this.
The specification also leaves a great deal to the deployer.
A deployment-pattern literature has grown around MCP to
record what~\cite{bridging_protocol_prod}. That work is a
practitioner account rather than a specification, so it
informs no rubric score. It is still where the gap between
a specified tool binding and a running one is documented.

\paragraph{ACP deployments.}
IBM's watsonx Orchestrate is the canonical ACP
deployment, at GA status. The BeeAI internal
incubator hosts additional pilot-stage deployments.
ACP's deployment footprint is smaller than MCP's or
A2A's, consistent with the rubric profile
(\S\ref{sec:industry:acp}): a structurally
sound tool-binding protocol that arrived later than
MCP into a market that had largely chosen MCP.

\paragraph{AGNTCY deployments.}
The Cisco-led AGNTCY effort has one named pilot
deployment at the cutoff (Webex AI Agent). The
operational pattern AGNTCY targets is transport-layer
integration of identity, routing, and observability
across heterogeneous agent stacks. That pattern is
structurally later-stage than the application-layer A2A
and MCP deployments. The cutoff reflects this maturity gap.

\paragraph{The composition pattern.}
A significant fraction of the surveyed enterprise SaaS
deployments run MCP \emph{inside} their A2A agents,
explicitly composing the tool-binding and peer-messaging
layers (\S\ref{sec:taxonomy:complementary}). Salesforce
Agentforce, for example, exposes A2A externally to other
agents and consumes MCP internally for tool integration.
The architectural argument that the two protocols
compose has produced a deployment reality where they
actually do.

\subsection{Agentic commerce and payments}
\label{sec:adoption:commerce}

The commerce vertical's deployment pattern is structurally
distinct from enterprise SaaS. Rather than dozens of
vendors each exposing their own agents, the cluster
consists of a small number of payment networks each
specifying their own settlement protocol. Status claims
follow the same Tier-4 sourcing
discipline~\cite{vendor_announcements_2026}; the four
settlement specifications themselves
(AP2~\cite{spec_ap2}, Visa TAP~\cite{spec_visa_tap},
Mastercard Agent Pay, ERC-8004~\cite{spec_erc8004}) carry
the behavioural claims about the
intent-mandate-receipt envelope. The four
candidates are AP2, Visa TAP, Mastercard Agent Pay, and
ERC-8004. They are competitors at the spec level but a
co-stack at the architectural level
(\S\ref{sec:emerging:commerce}). UCP~\cite{spec_ucp} sits
above all four and is treated below on a different basis,
because its adoption evidence at the cutoff is institutional
rather than deployment-side.

\paragraph{PayPal AP2.}
PayPal's reference implementation of AP2 is the canonical
shipping deployment of the protocol at GA, integrated
into PayPal's existing payment flows. The deployment
demonstrates the intent-mandate-receipt envelope in
production: an agent acting on behalf of a principal
proposes a transaction, the principal's mandate is
verified against the proposed intent, and the receipt of
execution is cryptographically attested through PayPal's
existing settlement chain.

\paragraph{Visa TAP.}
Visa Trusted Agent Protocol is at GA status with select
issuer-bank deployments. The Visa specification binds the
intent-mandate-receipt envelope to the EMVCo card-network
settlement substrate; the deployment pattern is that
issuer banks integrate TAP-attested agent transactions
into their existing card-network flows rather than
running a parallel rail. The architectural choice has
operational consequences: TAP-enabled transactions
flow through the same authorisation and clearing
infrastructure as card-present transactions, which makes
deployment fast but couples TAP to the existing
card-network governance.

\paragraph{Mastercard Agent Pay.}
Mastercard's analogue is in pilot status, targeting a
similar issuer-bank integration model. The cross-network
interoperability question
(\S\ref{sec:emerging:commerce}) asks whether a single
agent mandate is portable across PayPal, Visa, and
Mastercard rails. Both Visa and Mastercard acknowledge
it as a priority. It has not been resolved at
the spec level at the cutoff. The AAIF Agentic Commerce
Working Group~\cite{lf_aaif} is the venue where the
harmonisation is being negotiated.

\paragraph{ERC-8004 reference implementations.}
The on-chain ERC-8004 effort~\cite{spec_erc8004} is at
pilot status with reference implementations on
Ethereum and EVM-compatible chains. The pattern is the
direct on-chain analogue of the off-chain commerce
protocols: the mandate is a smart contract, the receipt
is a transaction record, and the principal-to-agent
binding is expressed as a signature against the
mandate contract. The pilot deployments are concentrated
in the agentic-DeFi sector: agents that execute DeFi
strategies on behalf of crypto-native principals. They
are a useful demonstration of the layered taxonomy
extending into substrates the off-chain payment networks
cannot reach.

\paragraph{UCP.}
UCP~\cite{spec_ucp} is the one commerce-layer protocol whose
adoption evidence at the cutoff is institutional rather than a
named deployment. It is founded by Google and Shopify, and its
Tech Council also seats Etsy, Target, and Wayfair
(\S\ref{sec:emerging:commerce}). Shopify's participation is
the notable part, because it puts merchant-side distribution
and a search-and-assistant platform in the same governing
body. That is a statement of commercial intent, not a
deployment claim. No merchant or platform implementation was
publicly verifiable at 2026-Q1, so the table records no status
for UCP. The survey's sourcing discipline
(\S\ref{sec:adoption}) does not let a membership list stand in
for a shipped integration. Whether the roster converts is the
question the next revision should answer.

\paragraph{Adjacent A2A and MCP deployments.}
Two adjacent deployments are worth flagging because they
demonstrate the composition pattern: Stripe's MCP
implementation (at GA) exposes Stripe's payment APIs as
MCP tools, allowing agents to invoke payments
imperatively; Shopify's A2A merchant-agent pilot exposes
merchant-side agents to upstream orchestrators over
A2A, with the actual payment transaction executed
through the commerce stack underneath. Neither
deployment is a commerce-settlement protocol in itself.
Both rely on AP2 or Visa TAP for the
intent-mandate-receipt envelope. But both illustrate
that the commerce layer composes with the peer-messaging
and tool-binding layers in production.

\subsection{Healthcare}
\label{sec:adoption:healthcare}

The healthcare vertical's deployment landscape is
sparser in public disclosure than the enterprise SaaS
or commerce verticals, for structural reasons that
the survey should name explicitly rather than working
around. Pilot announcements
(\cite{vendor_announcements_2026}) and the proprietary
intra-organisational deployments noted below are the
two visible patterns at the 2026-Q1 cutoff.

\paragraph{The disclosure gap.}
HIPAA and equivalent regulations create incentives for
healthcare vendors to integrate inter-agent
capabilities behind proprietary interfaces rather than
through public protocol specifications. The visible
result is that healthcare deployments that exist are
often invisible at the protocol level even when they
exist in production. Hippocratic AI's clinical agents,
for example, are deployed at GA in multiple hospital
systems but use proprietary internal interfaces; the
public claim is ``LLM agents'' rather than ``MCP
servers'' or ``A2A agents.''

\paragraph{Named A2A pilots.}
Epic agent connectivity (in pilot) and the Cedars-Sinai
plus Anthropic collaboration (in pilot) are the two
named A2A-relevant deployments at the cutoff. Both
target the multi-agent coordination pattern in clinical
workflows: a primary care agent invoking specialist
agents, a diagnostic agent invoking imaging-analysis
agents. The pilots use A2A's AgentCard convention for
discovery but layer their own authorisation policies
on top to satisfy HIPAA's principal-delegation
requirements. The rubric scores that layering low
on Axis 3 (\S\ref{sec:rubric:axis3}) precisely because
A2A does not provide it natively.

\paragraph{Named MCP deployments.}
Nuance DAX Copilot (now part of Microsoft) operates at
GA status and uses MCP-style tool integration
internally. Hippocratic AI's clinical agent platform
exposes MCP tools for selected clinical workflows.
Both are GA but proprietary in important ways: the
MCP servers they expose are not third-party
addressable. The deployment pattern in healthcare
favors what we will call \emph{intra-organisational
MCP}: MCP as an architecture choice within a single
vendor's stack, not as an interoperability protocol
across vendors.

\paragraph{The compliance-shaped roadmap.}
Healthcare adoption of A2A and MCP at the
cross-organisational level requires a delegation layer
that the contemporary protocols do not provide at
spec level (\S\ref{sec:security:authz}, O1). The HL7
FHIR community has begun discussion of an
inter-agent extension to FHIR-based identity flows
that would carry HIPAA-compliant delegation tokens
alongside A2A messages; this work is at the
pre-specification stage at the cutoff. Until it
matures, the healthcare adoption pattern is likely to
remain intra-organisational rather than
inter-organisational, with the
disclosure-gap consequence that the publicly visible
deployment list under-represents what is actually
running.

\subsection{Finance}
\label{sec:adoption:finance}

The finance vertical has its own version of the
disclosure problem: financial-services vendors deploy
inter-agent systems but seldom publish protocol
specifications, both for competitive reasons and
because financial regulation favours proprietary
audit-and-control surfaces over standard ones. The
named pilots below are the visible public
deployments~\cite{vendor_announcements_2026}; the
custom-protocol pattern that dominates the vertical is
discussed in the closing paragraph.

\paragraph{Named A2A and MCP deployments.}
Bloomberg AI's selected-feeds pilot uses A2A for
inter-agent coordination across Bloomberg's data
products. BlackRock Aladdin AI's tool integration
layer is in pilot with MCP. JPMorgan's IndexGPT
family of agents is in pilot and uses MCP internally
for tool calls. All three are pilot rather than GA;
all three target intra-organisational coordination
rather than cross-organisational interoperability.

\paragraph{The settlement-adjacent commerce angle.}
Finance adoption of the agentic commerce stack is
concentrated in settlement-adjacent use cases. DTCC
has pilot work on settlement-agent flows that use the
commerce-settlement layer's primitives (the
intent-mandate-receipt envelope) for cryptographically
attested securities settlement. The pattern is the
finance analogue of the PayPal AP2 deployment but in
a different settlement substrate. The pilot is the
clearest indication at the cutoff that the
commerce-settlement layer's primitives generalise
beyond consumer payments into institutional finance.

\paragraph{Custom protocols dominate.}
The dominant finance pattern at the cutoff remains
custom protocols (FIX-adjacent extensions, proprietary
banking APIs, internal multi-agent frameworks) rather
than the standardised inter-agent protocols that this
survey scores. The reason is not technical
inadequacy of A2A or MCP. They would serve most
finance-internal use cases adequately. The reason is a
regulatory and operational preference for proprietary
interfaces with proprietary audit chains. The
cross-organisational financial-data pattern (a buy-side
agent invoking a sell-side agent across firms) is the
case where standardised protocols would matter most;
no major public deployment exists at the cutoff. The
governance analysis in \S\ref{sec:governance} returns
to this question as an institutional matter.

\subsection{What the adoption pattern says about the protocols}
\label{sec:adoption:synthesis}

Three synthesising observations close the section.

\paragraph{Adoption clusters by layer, not by vendor.}
The tool-binding layer (MCP, ACP) and the peer-messaging
layer (A2A) cluster differently across verticals. MCP
adoption is densest in developer tooling, where the
tool-binding pattern is the natural fit. A2A adoption is
densest in business applications, where the peer-messaging
pattern is the natural fit. The commerce-settlement
layer is concentrated in the payment networks, where the
intent-mandate-receipt envelope is the natural fit. The
clusters indicate that the three-layer taxonomy
(\S\ref{sec:taxonomy:layers}) maps to real
deployment-shape decisions rather than to a purely
analytical convenience.

\paragraph{Cross-organisational deployment lags
intra-organisational.}
Every vertical exhibits the same pattern: an organisation
adopts an inter-agent protocol within its own boundaries
faster than it integrates across organisational
boundaries. The schema-fragmentation risk
(\S\ref{sec:semantic:fragmentation}) and the
authorisation-propagation gap (\S\ref{sec:security:authz})
are the two technical causes the
survey has identified for this lag; the
non-technical causes (competitive incentives,
disclosure preferences, regulatory uncertainty) compound
them. The result is that the inter-agent protocols of
the 2024 to 2026 wave are, at the cutoff, principally
intra-organisational integration substrates, with
cross-organisational adoption as the next-cycle question.

\paragraph{Adoption density is uncorrelated with rubric
strength.}
The protocols with the highest rubric coverage
(ACNBP at 10/12 on Axis 8, LOKA at level 3 on Axis 9)
have the smallest deployment footprints; the protocols
with the lowest rubric coverage (A2A, MCP, and ACP, all
at 1/12 on Axis 8) have the largest. The pattern is consistent with
the institutional analysis of \S\ref{sec:governance}:
governance maturity, vendor backing, and
ecosystem support drive adoption more than technical
rigour. The rubric measures technical quality at the
specification level; the adoption table measures
institutional traction; the two are
correlated only loosely and the survey should not pretend
otherwise.
